% =========================================================
% Compléments M1 — Brownien (B. Théorie)
% =========================================================

\subsection*{(1) Ce que la définition implique vraiment : structure gaussienne et covariance}
La définition d'un mouvement brownien standard $(W_t)_{t\ge 0}$ peut se résumer en une phrase :
\emph{processus gaussien centré, à trajectoires continues, de covariance $\Cov(W_s,W_t)=\min(s,t)$}.
En effet :
\begin{itemize}[leftmargin=*]
\item Les incréments stationnaires gaussiens impliquent que, pour tout $t$, $W_t\sim\Normal(0,t)$.
\item Pour $0\le s\le t$, on écrit $W_t=W_s+(W_t-W_s)$ avec $(W_t-W_s)$ indépendant de $\F_s$ et centré. Alors
\[
\Cov(W_s,W_t)=\E[W_sW_t]=\E\!\big[W_s(W_s+(W_t-W_s))\big]=\E[W_s^2]=s.
\]
\end{itemize}
La covariance $\min(s,t)$ est la façon la plus compacte de retenir la dépendance temporelle : \emph{le passé jusqu'à $\min(s,t)$ est commun, puis les incréments au-delà sont indépendants}.

\subsection*{(2) Propriété de Markov et propriété de martingale (formalisées)}
\paragraph{Markov.}
Le Brownien est un processus de Markov : conditionnellement à $W_t$, le futur $(W_{t+u}-W_t)_{u\ge 0}$ est indépendant du passé et suit une loi brownienne.
Concrètement, pour une fonction mesurable bornée $g$,
\[
\E[g(W_{t+u})\mid \F_t]=\E[g(W_t+(W_{t+u}-W_t))\mid \F_t]
=h_u(W_t),
\]
où $h_u(x)=\E[g(x+Z)]$ avec $Z\sim\Normal(0,u)$.

\paragraph{Martingale.}
$(W_t)$ est une martingale sous $\Pmeas$ :
\[
\E[W_t\mid \F_s]=W_s,\qquad 0\le s\le t.
\]
\emph{Preuve courte.} Écrire $W_t=W_s+(W_t-W_s)$ et utiliser $\E[W_t-W_s\mid \F_s]=0$ par indépendance et centrage des incréments.

\paragraph{Pourquoi ça compte en finance ?}
Deux raisons structurantes :
\begin{itemize}[leftmargin=*]
\item La martingale capture l'idée «pas de dérive prédictible une fois l'information prise en compte». Sous $\Qmeas$, ce seront les \emph{prix actualisés} qui deviennent des martingales.
\item Le Markov est le pont vers les PDE : si le prix $V(t,S_t)$ dépend seulement de $(t,S_t)$ (pas du chemin), alors $V$ peut souvent être caractérisé par une équation de Kolmogorov / une PDE.
\end{itemize}

\subsection*{(3) Invariance d'échelle (scaling) : d'où vient le $\sqrt{\Delta t}$}
La propriété de scaling dit :
\[
(W_{ct})_{t\ge 0}\ \overset{d}{=}\ (\sqrt{c}\,W_t)_{t\ge 0}.
\]
\emph{Lecture.} Quand on change d'unité de temps (par exemple jours $\to$ années), l'échelle du bruit change comme $\sqrt{c}$.
C'est la justification mathématique de la règle opérationnelle :
\[
\Delta W \approx \sqrt{\Delta t}\,Z,\qquad Z\sim\Normal(0,1).
\]
\emph{Preuve (sur les lois finies-dimensionnelles).} Pour des temps $0\le t_1<\dots<t_n$, le vecteur $(W_{ct_1},\dots,W_{ct_n})$ est gaussien centré, de covariance
\[
\Cov(W_{ct_i},W_{ct_j})=\min(ct_i,ct_j)=c\,\min(t_i,t_j),
\]
qui est exactement la covariance de $(\sqrt{c}\,W_{t_1},\dots,\sqrt{c}\,W_{t_n})$.

\subsection*{(4) Propriétés de trajectoire : continuité oui, dérivabilité non}
Deux faits contre-intuitifs mais essentiels :
\begin{itemize}[leftmargin=*]
\item Les trajectoires de $W$ sont continues, mais \textbf{presque sûrement nulle part dérivables}.
\item La \textbf{variation totale} de $W$ sur tout intervalle est infinie (presque sûrement).
\end{itemize}
\paragraph{Pourquoi c'est important ?}
Parce que cela explique pourquoi le calcul différentiel «classique» échoue : on ne peut pas traiter $\dd W_t$ comme une différentielle ordinaire, et les termes d'ordre 2 (variation quadratique) deviennent déterminants.

\paragraph{Hölder (ordre de grandeur précis).}
On peut montrer que pour tout $\alpha<\tfrac12$, $W$ est Hölder-$\alpha$ sur $[0,T]$ p.s., et pour $\alpha\ge\tfrac12$ ce n'est plus vrai.
Intuition : $\Delta W=O(\sqrt{\Delta t})$ est exactement la frontière entre «variation finie» et «variation infinie».

\subsection*{(5) Principe de réflexion : maximum, barrières, temps de premier passage}
Le principe de réflexion est l'outil classique pour obtenir des lois de maximum/minimum du Brownien.
Soit $M_T=\sup_{0\le t\le T} W_t$.
Alors, pour $a>0$ :
\[
\Prob(M_T\ge a)=2\Prob(W_T\ge a)=2\bigl(1-N(a/\sqrt{T})\bigr).
\]
\paragraph{Idée de preuve (niveau M1).}
On considère l'événement où le chemin franchit $a$ avant $T$, et on «réfléchit» la trajectoire après le premier passage $\tau_a=\inf\{t:W_t=a\}$.
Cette transformation envoie bijectivement les trajectoires qui franchissent $a$ et finissent \emph{sous} $a$ vers des trajectoires qui ne franchissent pas mais finissent \emph{au-dessus} de $a$.
On obtient alors une relation de symétrie sur les probabilités, qui se traduit en gaussien par la formule ci-dessus.

\paragraph{Lien avec finance.}
Pour une barrière (knock-out/knock-in), la probabilité de franchissement d'un niveau intervient directement dans le prix.
En Black--Scholes, on transporte ce type de résultat via la transformation logarithmique du GBM.

\subsection*{(6) Intégrale d'Itô : construction rigoureuse en $L^2$ (version cours)}
On veut définir $\int_0^T H_t\,\dd W_t$ pour des processus adaptés $H$.
La construction standard (niveau M1) suit trois étapes.

\paragraph{Étape 1 : processus simples (prévisibles en escalier).}
On appelle \emph{processus simple prévisible} sur $[0,T]$ toute forme
\[
H_t=\sum_{k=0}^{n-1} H_{t_k}\,\1_{(t_k,t_{k+1}]}(t),
\qquad 0=t_0<\dots<t_n=T,
\]
où chaque $H_{t_k}$ est $\F_{t_k}$-mesurable (donc «décidé» au début de l'intervalle).
On définit alors
\[
\int_0^T H_t\,\dd W_t\ \triangleq\ \sum_{k=0}^{n-1} H_{t_k}\bigl(W_{t_{k+1}}-W_{t_k}\bigr).
\]

\paragraph{Étape 2 : isométrie d'Itô (preuve sur les simples).}
Pour un tel $H$, on a
\[
\E\!\left[\left(\int_0^T H_t\,\dd W_t\right)^2\right]
=\E\!\left[\int_0^T H_t^2\,\dd t\right].
\]
\emph{Preuve (calcul explicite).}
Écrire $I=\sum_k H_{t_k}\Delta W_k$ avec $\Delta W_k=W_{t_{k+1}}-W_{t_k}$.
Alors
\[
\E[I^2]=\E\!\left[\sum_{k,\ell} H_{t_k}H_{t_\ell}\Delta W_k\Delta W_\ell\right].
\]
Pour $k\neq \ell$, les incréments sont indépendants, centrés, et l'un des deux est indépendant de l'information contenant l'autre coefficient : les termes croisés ont espérance nulle.
Donc
\[
\E[I^2]=\sum_k \E\!\left[H_{t_k}^2\,\E[(\Delta W_k)^2\mid \F_{t_k}]\right]
=\sum_k \E[H_{t_k}^2]\,(t_{k+1}-t_k)
=\E\!\left[\int_0^T H_t^2\,\dd t\right].
\]

\paragraph{Étape 3 : complétion en $L^2$.}
On définit l'espace
\[
\mathcal{H}^2=\left\{H\ \text{prévisible}:\ \E\left[\int_0^T H_t^2\,\dd t\right]<\infty\right\}.
\]
Les processus simples sont denses dans $\mathcal{H}^2$ pour la norme $\|H\|_{2}^2=\E[\int_0^T H_t^2\dd t]$.
On choisit une suite $(H^n)$ simple qui converge vers $H$ dans cette norme, et on pose
\[
\int_0^T H_t\,\dd W_t \ \triangleq\ \lim_{n\to\infty} \int_0^T H^n_t\,\dd W_t
\quad \text{dans } L^2.
\]
L'isométrie garantit que la limite existe et ne dépend pas de l'approximation choisie.

\paragraph{Conséquence : martingale et variance.}
Pour $H\in\mathcal{H}^2$, le processus
\[
M_t=\int_0^t H_s\,\dd W_s
\]
est une martingale, et son «énergie» est
\[
\E[M_T^2]=\E\!\left[\int_0^T H_t^2\,\dd t\right].
\]
En finance, cette identité est la brique qui relie un contrôle quadratique (risque) aux expositions instantanées $H_t$.

\subsection*{(7) Itô vs Stratonovich (en une page) : convention et sens économique}
Deux conventions d'intégration apparaissent souvent :
\begin{itemize}[leftmargin=*]
\item Itô : on gèle l'intégrande au début de l'intervalle $(t_k,t_{k+1}]$.
\item Stratonovich : on gèle «au milieu» (ou plus exactement on prend une limite symétrique).
\end{itemize}
En finance, Itô est naturel pour un raisonnement \emph{non anticipatif} : une stratégie de trading ne peut pas dépendre de l'incrément futur $\Delta W_k$ sur l'intervalle.
Stratonovich est plus fréquent en physique (modélisation de bruits «lisses» approchés).
La conversion (en dimension 1) est :
\[
\int_0^T H_t\circ \dd W_t
=\int_0^T H_t\,\dd W_t+\frac12\int_0^T \frac{\partial H_t}{\partial x}\,\dd t
\quad \text{(si }H_t=H(t,X_t)\text{)}.
\]
Le terme $\tfrac12$ est la «correction» qui renvoie au chapitre Variation quadratique.

\subsection*{(8) Brownien multi-dimensionnel et corrélations (utile pour Heston/SABR)}
Un brownien $d$-dimensionnel $\bm{W}_t=(W_t^{(1)},\dots,W_t^{(d)})$ a des composantes browniens et des covariations
\[
[W^{(i)},W^{(j)}]_t=\rho_{ij}t,
\]
où la matrice $\rho$ (symétrique, définie positive) code les corrélations instantanées.
Une représentation standard est :
\[
\bm{W}_t = L\bm{B}_t,
\]
où $\bm{B}$ est un brownien à composantes indépendantes et $L$ une matrice telle que $LL^\top=\rho$ (décomposition de Cholesky).
Cette écriture est cruciale dès qu'on couple prix et variance (Heston) ou prix et volatilité (SABR).
